\chapter{Testtext}
\section{Theoretische Grundlagen}
\subsection{Durchflussmessung}

Die Durchflussmessung in prozesstechnischen Anlagen ist ein komplexes Thema, das von einer Vielzahl an Parametern und Rahmenbedingungen abhängt. In Rohrleitungen wird der Durchfluss maßgeblich durchFaktoren wie die spezifischen Stoffeigenschaften des Fluids, die Charakteristiken der Fördermaschine, vorhandene Höhen- und Druckunterschiede sowie die Eigenschaften der Rohrleitung selbst, wie Länge oder Rauigkeit, beeinflusst.\\
Integrierte Durchflussstelleinrichtungen und Messsysteme können den Durchfluss mehr oder weniger stark modifizieren. Die Vielfalt der möglichen Prozessmedien erfordert entsprechend differenzierte Messmethoden. Diese müssen die chemisch-physikalischen Eigenschaften wie Zusammensetzung, Leitfähigkeit und Viskosität ebenso berücksichtigen wie mögliche Aggregatzustände und Gemischzusammensetzungen. Falls die Zuverlässigkeit im Vordergrund steht, kann ein Durchflusswächter, der ein binäres Signal liefert, ausreichend sein, doch meistens ist die Genauigkeit ein wichtiges Auswahlkriterium.

\subsection{Messverfahren}
\subsubsection{Wirkdruckverfahren}

Das Wirkdruckverfahren basiert auf der Bernoulli-Gleichung~(\ref{eqn: Bernoulli}) und der Kontinuitätsgleichung~(\ref{eqn: Konti}).
\begin{align}
	\label{eqn: Bernoulli}
	p_1 + \frac{1}{2} \rho v_1^2 + \rho g h_1 &= p_2 + \frac{1}{2} \rho v_2^2 + \rho g h_2\\
	\label{eqn: Konti}
	A_1 v_1 &= A_2 v_2
\end{align}
\begin{enumerate}[]
	\item[] $p$: Druck
	\item[] $\rho$: Dichte
	\item[] $v$: Geschwindigkeit
	\item[] $g$: Erdbeschleunigung
	\item[] $h$: Höhe
	\item[] $A$: Querschnittsfläche
\end{enumerate}

Bei der Wirkdruckmessung wird ein Strömungswiderstand, i.d.R. Düsen oder Blenden nach DIN 1952, in die Strömung eingebaut, der einen Druckunterschied (Wirkdruck) zwischen beiden Seiten verursacht. Der nichtlineare Zusammenhang zwischen Volumenstrom $Q_v$ und Wirkdruck $\Delta$p wird bei Vernachlässigung der Reibung durch Gleichung~() beschrieben.